\documentclass[../main.tex]{subfiles}

\begin{document}

    \section{Code}\label{sec:code}

    The code for the report has been condensed into one executable file.
    The original GitHub project is linked \href{https://github.com/TheSilverLeopard/University/tree/projects/MTH2005%20Modelling%20-%20Theory%20and%20Practice/Coursework%201}
    {here}.

    \subsection*{Usage of modules}\label{subsec:usage-of-modules}
    The code has been structured in a way that allows ease of use and reusability. \\
    \\
    Below is a file called \texttt{main.py}, it is an example use of the \texttt{case\_study} module and a shorthand for
    running any code for the questions:
    \lstinputlisting[label={lst:code-example-usage}, style=python]{../code-examples/usage-example.py}

    \subsubsection{Fixed-point method}\label{subsubsec:fixed-point-method-appendix}
    Running this in your console will execute the fixed-point method code and
    then save the plots to the output directory provided.
    \begin{lstlisting}[style=console,label={lst:appendix-fixed-point-usage}]
python main.py fixed-point <output-directory>
    \end{lstlisting}
    \\
    For example:
    \begin{lstlisting}[style=console,label={lst:appendix-fixed-point-example}]
C:\User\CaseStudy> python main.py fixed-point ./images
Root: 1.2345
Error: 0.1234
    \end{lstlisting}

    \subsubsection{2D Newton-Raphson}\label{subsubsec:2d-nr-appendix}
    Running this in your console will execute the 2D Newton-Raphson code.
    \begin{lstlisting}[style=console,label={lst:appendix-2d-nr-usage}]
python main.py 2d-nr
    \end{lstlisting}
    \\
    For example:
    \begin{lstlisting}[style=console,label={lst:appendix-2d-nr-example}]
C:\User\CaseStudy> python main.py 2d-nr
Root: [-1.2345, 1.2345]
Root verification: [0. 0.]
Hessian eigenvalues: [-12, 34]
Stationary point nature: saddle point
    \end{lstlisting}

    \subsubsection{Boundary-value matrix problem}\label{subsubsec:bvp-appendix}
    Running this in your console will execute the Boundary-value matrix problem code and
    then save the plots to the output directory provided.
    \begin{lstlisting}[style=console,label={lst:appendix-bvp-usage}]
python main.py boundary <output-directory>
    \end{lstlisting}
    \\
    For example:
    \begin{lstlisting}[style=console,label={lst:appendix-bvp-example}]
C:\User\CaseStudy> python main.py boundary ./images
Order: 1.00
    \end{lstlisting}

    \subsubsection{Condensing the module}
    Running this in your console will take the input directory, condense the module into a single file,
    and save it to the output directory provided.
    \begin{lstlisting}[style=console,label={lst:appendix-condensing-module-usage}]
python main.py condense <input-directory> <output-directory>
    \end{lstlisting}
    \\
    For example:
    \begin{lstlisting}[style=console,label={lst:appendix-condensing-module-example}]
C:\User\CaseStudy> python main.py condense ./case_study ./output.py
    \end{lstlisting}

    \subsection*{Condensed module code}\label{subsec:condensed-code}
    I would put the following code into \texttt{case\_study.py} for the \texttt{main.py} code from above to work
    directly with this.
    Or you could remove the line that imports \texttt{case\_study.py} and paste the contents of
    \textt{main.py} directly under the condensed code.
    \lstinputlisting[label={lst:code_appendix}, style=python]{../code-examples/output.py}

\end{document}